% C.5 AB∥CD，EF 截于 G,H；∠EGB=65°
\documentclass[tikz,border=4pt]{standalone}
\usepackage{ctex}
\usepackage{amsmath}
\usetikzlibrary{calc, arrows.meta}
\begin{document}
\begin{tikzpicture}[scale=1.0, thick]
  % AB horizontal at y=2; CD at y=0
  \draw[-{Stealth[length=3pt]}] (-3,2) -- (3.2,2) node[right] {$B$};
  \node[left] at (-3,2) {$A$};
  \draw[-{Stealth[length=3pt]}] (-3,0) -- (3.2,0) node[right] {$D$};
  \node[left] at (-3,0) {$C$};
  % EF: line through G=(0,2) and H=(... ) making 65° with GB (i.e. with +x dir)
  % slope = tan(65°+90°)? Actually ∠EGB=65 between ray GE and ray GB.
  % Let E above, F below. GE makes angle 180-65=115 from +x.
  \pgfmathsetmacro{\ang}{115}
  \coordinate (G) at (0,2);
  \coordinate (H) at ({0 + 2/tan(\ang-90)*(-1)}, 0); % compute intersection: line from G in direction (cos\ang, sin\ang); reaches y=0 when t*sin\ang=-2.
  % simpler: use slope. Direction: (cos115, sin115). From G(0,2): param t. y=2+t*sin115=0 => t=-2/sin115. x= t*cos115.
  \pgfmathsetmacro{\tval}{-2/sin(\ang)}
  \pgfmathsetmacro{\Hx}{\tval*cos(\ang)}
  \coordinate (H) at (\Hx, 0);
  % E is extension above G; F extension below H
  \pgfmathsetmacro{\Ex}{1.2*cos(\ang)}
  \pgfmathsetmacro{\Ey}{2 + 1.2*sin(\ang)}
  \coordinate (E) at (\Ex,\Ey);
  \pgfmathsetmacro{\Fx}{\Hx - 1.2*cos(\ang)}
  \pgfmathsetmacro{\Fy}{- 1.2*sin(\ang)}
  \coordinate (F) at (\Fx,\Fy);
  \draw[-{Stealth[length=3pt]}] (F) -- (E);
  \filldraw (G) circle (1.2pt) node[above right] {$G$};
  \filldraw (H) circle (1.2pt) node[below right] {$H$};
  \node[above] at (E) {$E$};
  \node[below] at (F) {$F$};
  % mark ∠EGB
  \draw (0.6,2) arc (0:\ang:0.6);
  \node[font=\footnotesize] at ({0.35*cos(\ang/2)}, {2 + 0.35*sin(\ang/2)}) {$65°$};
\end{tikzpicture}
\end{document}
